\subsection{房与单房}

用 \(t_{r}\) 表示 t 中使 exp x 为正则的元素 \(x \in t\) 所成的集合，换言之，属于 \(T_{r}\) 的元素所成的集合。t 的元素 x 属于 \(t - t_{r}\) 当且仅当存在根 \(\alpha \in \mathrm{R}(\mathrm{G}, \mathrm{T})\) 使得 \(\delta(\alpha)(x) \in 2\pi i\mathbf{Z}\)。对每个根 \(\alpha \in \mathrm{R}(\mathrm{G}, \mathrm{T})\) 与每个整数 n，用 \(H_{\alpha,n}\) 表示诸 \(x \in t\) 中满足 \(\delta(\alpha)(x) = 2\pi in\) 者所成的集合。诸 \(H_{\alpha,n}\) 称为 t 的奇异超平面，而 \(t - t_{r}\) 是诸奇异超平面之并。t 的单房是 \(t_{r}\) 的连通分支，而房是 t 中过原点的那些奇异超平面（即 \(H_{\alpha,0} = \operatorname{Ker}\delta(\alpha), \alpha \in \mathrm{R}(\mathrm{G}, \mathrm{T})\)）之并的补集的连通分支。

我们有 \(\varGamma(\mathrm{T})\subset\mathfrak{t}-\mathfrak{t}_{r}\)；用 N(G,T) 表示由结点向量生成的 \(\varGamma(\mathrm{T})\) 的子群（§4，第 5 节）；由 §4 第 6 节命题 11，商 \(\varGamma(\mathrm{T})/\mathrm{N}(\mathrm{G},\mathrm{T})\) 可等同于 G 的基本群。

最后，用 W 表示 G 关于 T 的 Weyl 群，把它视为 T 与 t 的自同构群，并用 \(W_{a}\)（相应地 \(W_{a}^{\prime}\)）表示由 W 与诸平移 \(t_{\gamma}: x \mapsto x + \gamma\)（\(\gamma \in \mathrm{N}(\mathrm{G}, \mathrm{T})\)，相应地 \(\gamma \in \Gamma(\mathrm{T})\)）所生成的仿射空间 t 的自同构群。

设 \(w \in W\)，\(\gamma \in \Gamma(T)\)，\(\alpha \in R(G, T)\) 且 \(n \in Z\)。我们有：

\[
w (\mathrm{H} _ {\alpha , n}) = \mathrm{H} _ {w \alpha , n}, t _ {\gamma} (\mathrm{H} _ {\alpha , n}) = \mathrm{H} _ {\alpha , n + \langle \gamma , \alpha \rangle}.
\]

由此可知，对所有房 C 与所有 \(w \in W\)，\(w(\mathrm{C})\) 是一个房，且对所有单房 A 与所有 \(w \in W_{a}^{\prime}\)，\(w(\mathrm{A})\) 是一个单房。注意，当 \(\mathrm{X}(\mathrm{T}) \otimes \mathbf{R}\) 等同于 \(t^{*}\)（借助同构 \((2\pi i)^{-1}\delta\)）时，t 的诸单房与群 \(W_{a}\) 就是相伴于根系 \(\mathrm{R}(\mathrm{G}, \mathrm{T})\) 的诸单房与仿射 Weyl 群（第六章，§2，第 1 节）。

命题 2. a) 群 \(W_{a}\)（相应地 \(W_{a}^{\prime}\)）是 W 与 N(G, T)（相应地 \(\Gamma(T)\)）的半直积；子群 \(W_{a}\)（\(W_{a}^{\prime}\) 的）是正规的。

\begin{enumerate}
\def\labelenumi{\alph{enumi})}
\setcounter{enumi}{1}
\item
  群 W（相应地 \(\mathrm{W}_a\)）单传递地作用于诸房（相应地诸单房）所成的集合上。
\item
  设 C 是一个房，A 是一个单房。则 \(\overline{C}\)（相应地 \(\overline{A}\)，相应地 A）是 W 在 t 上的作用（相应地 \(W_{a}\) 在 t 上的作用，相应地 \(W_{a}\) 在 \(t-t_{r}\) 上的作用）的基本域。若 \(x \in t_{r}\) 与 \(w \in W_{a}\) 满足 \(w(x) = x\)，则 w = Id。
\item
  对每个房 C，存在唯一的单房 A 使得 \(A \subset C\) 且 \(0 \in \overline{A}\)。对每个单房 A，存在唯一的 \(\gamma \in \mathrm{N}(\mathrm{G}, \mathrm{T})\) 使得 \(\gamma \in \overline{A}\)。
\end{enumerate}

若 \(w \in W\) 且 \(\gamma \in \Gamma(T)\)，则 \(wt_{\gamma}w^{-1} = t_{w(\gamma)}\)，且 \(wt_{\gamma}w^{-1}t_{\gamma}^{-1} = t_{w(\gamma)-\gamma}\)，其中 \(w(\gamma) - \gamma \in \mathrm{N}(\mathrm{G}, \mathrm{T})\)；由此立得 a)。命题的其余部分由第六章 §1 第 5 节与 §2 第 1、2 节得出。

推论 1. 设 A 是 t 的一个单房，\(\overline{A}\) 是它的闭包，\(H_{A}\) 是 A 在 \(W_{a}^{\prime}\) 中的稳定子群。

\begin{enumerate}
\def\labelenumi{\alph{enumi})}
\item
  群 \(W_{a}^{\prime}\) 是 \(H_{A}\) 与 \(W_{a}\) 的半直积。
\item
  指数映射 \(\overline{\mathrm{A}}\to \mathrm{T}\) 与典则单射 \(\mathrm{T}\rightarrow \mathrm{G}\) 通过过渡到商与子集，诱导出同胚
\end{enumerate}

\[
\begin{array}{l} \overline {{\mathrm{A}}} / \mathrm {H_ {A}} \to \mathrm{T} / \mathrm{W} \to \mathrm{G} / \mathrm{Int} (\mathrm{G}) \\ \mathrm{A} / \mathrm {H_ {A}} \to \mathrm {T_ {r}} / \mathrm{W} \to \mathrm {G_ {r}} / \mathrm{Int} (\mathrm{G}). \end{array}
\]

设 \(w' \in W_a'\)；则 \(w'(A)\) 是 \(t\) 的一个单房，且存在（命题 2 b)）唯一的 \(w\) 属于 \(W_a\)，使得 \(w(A) = w'(A)\)，即 \(w^{-1}w' \in H_A\)。由于 \(W_a\) 在 \(W_a'\) 中正规，这就证明了 \(a\))。

\(\overline{A}\) 到 t 中的典则单射诱导出一个连续双射 \(\theta: \overline{A} \to t/W_{a}\)（命题 2 c)），而它是同胚，因为 \(\overline{A}\) 是紧的。由于 \(W_{a}\) 在 \(W_{a}^{\prime}\) 中正规，群 \(H_{A}\) 典则地作用于 \(t/W_{a}\) 上（《代数》，第一章，§5，第 4 节），且 \(t/W_{a}^{\prime}\) 可等同于商 \((t/W_{a})/H_{A}\)；映射 \(\theta\) 与 \(H_{A}\) 的作用相容，故通过过渡到商诱导出同胚 \(\overline{A}/H_{A} \to t/W_{a}^{\prime}\)。进而，\(\exp_{\Gamma}\) 诱导出从 \(t/\Gamma(T)\) 到 T 的一个同胚，从而也诱导出从 \(t/W_{a}^{\prime}\) 到 T/W 的一个同胚。断言 b) 由这一点以及 §2 第 4 节命题 5 的推论 1 得出。

注. 1) 群 \(H_{A}\) 可自然地等同于 \(\Gamma(\mathrm{T})/\mathrm{N}(\mathrm{G},\mathrm{T})\)，从而也可等同于 \(\pi_{1}(\mathrm{G})\)。于是，当 G 单连通时，它就化为单位元素。

\begin{enumerate}
\def\labelenumi{\arabic{enumi})}
\setcounter{enumi}{1}
\item
  设 \(x \in \mathrm{A}\)；则 \(\exp x \in \mathrm{T}_r\)，故 \(\mathrm{Z}(\exp x)_0 = \mathrm{T}\)。进而，\(\exp x\) 是强正则的（第 1 节，注）当且仅当 \(w(x) \neq x\) 对所有异于恒等映射的 \(w \in \mathrm{W}_a'\) 成立。由推论 1，这也意味着 \(h(x) \neq x\) 对所有异于恒等映射的 \(h \in \mathrm{H}_{\mathrm{A}}\) 成立。特别地，若 \(\mathrm{G}\) 单连通，则 \(\mathrm{Z}_{\mathrm{G}}(t) = \mathrm{T}\) 对所有 \(t \in \mathrm{T}_r\) 成立，且 \(\mathrm{G}\) 的每个正则元都是强正则的。
\item
  \(W_{a}\) 的特殊点（第六章，§2，第 2 节）是 t 中满足 \(\delta(\alpha)(x) \in 2\pi i\mathbf{Z}\)（对所有 \(\alpha \in \mathrm{R}(\mathrm{G}, \mathrm{T})\)）的元素 x（同前处，命题 3），也就是满足 \(\exp x \in \mathrm{C}(\mathrm{G})\) 的元素 x（§4，第 4 节，命题 8）。对这样的元素 x，我们有 \(wx - x \in \mathrm{N}(\mathrm{G}, \mathrm{T})\)（对所有 \(w \in W\)）（第六章，§1，第 9 节，命题 27），故 x 在 \(W_{a}\) 与 \(W_{a}^{\prime}\) 中的稳定子群一致。设 S 是 \(\overline{A}\) 的特殊点所成的集合；由前述结果与推论 1 可知，群 \(H_{A}\) 自由地作用于 S 上，且指数映射诱导出从 \(S/H_{A}\) 到 \(C(G)\) 的一个双射。
\end{enumerate}

推论 2. 设 C 是 t 的一个房，\(\overline{C}\) 是它的闭包。典则单射 \(\overline{C} \rightarrow t \rightarrow g\) 通过过渡到商诱导出同胚

\[
\overline {{\mathrm{C}}} \rightarrow \mathfrak {t} / \mathrm{W} \rightarrow \mathfrak {g} / \operatorname{Ad} (\mathrm{G}).
\]

典则映射 \(\overline{C} \rightarrow t\) 与 \(t \rightarrow t/W\) 是真的（《一般拓扑》，第三章，§4，第 1 节，命题 2 c)）。映射 \(\overline{C} \rightarrow t/W\) 连续、真且双射（命题 2 c)）；因而是同胚，于是结合 §2 第 5 节命题 6 的推论即得推论。

注 4. 用 \(\mathfrak{g}_{\mathrm{reg}}\) 表示 \(\mathfrak{g}\) 的正则元所成的集合（第七章，§2，第 2 节，定义 2），并令 \(\mathfrak{t}_{\mathrm{reg}} = \mathfrak{t} \cap \mathfrak{g}_{\mathrm{reg}}\)。对 \(x \in \mathfrak{t}\)，我们有

\[
\det (\mathrm{X} - \operatorname{ad} _ {\mathfrak {g}} x) = \mathrm{X} ^ {\dim t} \prod_ {\alpha \in \mathrm{R(G,T)}} (\mathrm{X} - \delta (\alpha) (x)),
\]

因而 \(t_{reg}\) 是 t 中满足 \(\delta(\alpha)(x) \neq 0\)（对所有 \(\alpha \in \mathrm{R}(\mathrm{G}, \mathrm{T})\)）的元素 x 所成的集合，也就是 t 的诸房之并（于是 \(t_r \subset t_{reg}\)）。由此，\(\overline{C} \cap t_{reg} = C\)，故我们有同胚

\[
\mathrm{C} \rightarrow \mathfrak {t} _ {\text { reg }} / \mathrm{W} \rightarrow \mathfrak {g} _ {\text { reg }} / \mathrm{Ad} (\mathrm{G}).
\]

推论 3. 假设 G 是单连通的；设 g 是 G 的一个正则元。存在 G 的一个极大环面 S 与 L(S) 的一个单房 A，二者均唯一确定，使得 \(g \in \exp(A)\) 且 \(0 \in \overline{A}\)。

可以假设 \(g\) 属于 \(\mathrm{T}_r\)（§2，第 2 节，定理 2）。设 \(x\) 是 \(\mathfrak{t}_r\) 中满足 \(\exp x = g\) 的一个元素，并设 \(\mathrm{A}'\) 是 \(\mathfrak{t}\) 中包含 \(x\) 的单房。诸单房 \(\mathrm{A}\)（\(\mathfrak{t}\) 中的）若满足 \(g \in \exp(\mathrm{A})\)，便是诸单房 \(\mathrm{A}' - \gamma\)（\(\gamma \in \Gamma(\mathrm{T})\)）；于是，断言由命题 2 d) 得出。

